% Generated draft: requires mathematical and source review.
\documentclass[11pt]{article}
\usepackage[margin=25mm]{geometry}
\usepackage{fontspec}
\setmainfont{Latin Modern Roman}
\usepackage{amsmath,amssymb,mathtools,unicode-math}
\setmathfont{Latin Modern Math}
\usepackage{enumitem,needspace,xurl,hyperref}
\setlength{\parindent}{0pt}
\setlength{\parskip}{5pt}
\setlength{\emergencystretch}{3em}
\newcommand{\N}{\mathbb{N}}
\newcommand{\Z}{\mathbb{Z}}
\newcommand{\Q}{\mathbb{Q}}
\newcommand{\R}{\mathbb{R}}
\newcommand{\C}{\mathbb{C}}
\newcommand{\F}{\mathbb{F}}
\newcommand{\A}{\mathbb{A}}
\newcommand{\PP}{\mathbb{P}}
\newcommand{\E}{\mathbb{E}}
\newcommand{\eps}{\varepsilon}
\newcommand{\dd}{\,\mathrm{d}}
\begin{document}
\section*{UnsolvedMath ID: 800865}
\textbf{Source problem code:} catalogue entry
\subsection{Definitions and mathematical statement}
A statistical model \(\mathcal M\) is a specified family of probability distributions on a common measurable sample space. An algebraic description of such a model uses polynomial equations in specified coordinates, possibly together with inequalities required for probabilities or model restrictions. The supplied source does not specify the sample space, coordinates, equations, inequalities, or class of models intended.

For a specified model \(\mathcal M\) and observed data \(d\) for which a likelihood function \(L_d\colon\mathcal M\to[0,\infty)\) has been defined, write
\[
\operatorname{MLE}_{\mathcal M}(d)
=\bigl\{P\in\mathcal M: L_d(P)\geq L_d(Q)\text{ for every }Q\in\mathcal M\bigr\}.
\]
This definition allows an empty set when the maximum is not attained and multiple elements when the maximizer is not unique. Defining \(L_d\) requires the sampling and likelihood conventions appropriate to the chosen model; these are not supplied.

\textbf{Research problem (Sturmfels, 2007, General Problem 4).} Study the geometry of maximum likelihood estimation for algebraic statistical models.

The notation above provides a conditional framework for interpreting maximum likelihood estimation, rather than a more specific conjecture attributed to the source. The source gives a research direction without specifying a geometric invariant, a proposed conclusion, a quantifier over models or data, or conventions governing boundary distributions. Consequently, it does not determine a unique theorem or conjecture to formalize.
\subsection{Short English statement}
Sturmfels asks for a study of the geometry of maximum likelihood estimation in algebraic statistical models. For a specified model and likelihood, the estimators are all distributions attaining the largest likelihood. The supplied problem leaves the model class, data assumptions, geometric target, and boundary conventions unspecified.
\subsection{Sources}
\begin{itemize}
\item Local UnsolvedMath registry, permanent ID \(800865\), exact source problem code \texttt{catalogue entry}, titled \emph{Open Problems in Algebraic Statistics --- General Problem 4}. The record attributes the problem to Sturmfels, with proposed year 2007.
\item Sturmfels, \emph{Open Problems in Algebraic Statistics} (2007), General Problem (\#4). Supplied source URL: \url{https://arxiv.org/abs/0707.4558}. The supplied provenance records access on 2026-10-16 and extraction by \texttt{arxiv\_latex}; these are supplied provenance claims, not independently verified access or source inspection.
\item The local registry labels the problem \texttt{open}. This is a supplied status label, not an independently verified assessment of the literature.
\end{itemize}
\end{document}
